\section{Executive summary}

The project investigates how to allocate a fixed reinforcement-learning budget
across user preferences when objectives conflict. A preference may be difficult,
uncertain, under-sampled, or currently learnable; these are different properties.
The research question is whether allocation signals predict future improvement
and improve performance relative to uniform and compute-matched random controls.

The preference-conditioned Q-learning pipeline, four allocation signals, paired
controls, fixed evaluation protocols, provenance tracking, and validation audits
have been implemented. The backbone implementation gate passed 103 tests;
the subsequent tooling and bounded-review gate passed 168 tests with formatting,
static analysis, security, and compilation checks. These are software checks,
not evidence of algorithm convergence or adaptive-allocation benefit.
Thirty primary course runs completed, but a validation audit found zero ore
returns throughout the Minecart block, invalidating a comparative interpretation.

Two predeclared backbone qualifications were then completed. V1 failed its
all-seeds gate; v2 tested proportional prioritized replay as a single change.
V2 source checks passed for all 30 checkpoints, but all three selected
checkpoints returned zero ore in both ore objectives across 160 episodes each.
\textbf{V2 therefore failed the unchanged qualification rule.} The current
research contribution is implementation and an audited failure diagnosis, not
evidence that adaptive preference allocation improves learning.

\section{Problem and research objectives}

\subsection{Why preference allocation matters}
In multi-objective reinforcement learning (MORL), the environment returns a
vector reward rather than one scalar. A preference vector $w$ specifies the
relative importance of objectives, with $w_i\geq0$ and $\sum_iw_i=1$.
A conditioned policy should adapt to these preferences without training a
separate policy for every trade-off. The unresolved training decision is where
to spend the next interaction or gradient update under a fixed total budget.

The course scope uses \texttt{minecart-v0} and
\texttt{mo-lunar-lander-v3}. Minecart is especially diagnostic because ore
collection and delivery require sparse, delayed successful behavior. A learner
that fails to acquire or retain this behavior cannot establish whether an
allocation signal is useful.

\subsection{Questions and bounded novelty}
The study asks:
\begin{enumerate}
\item Which signal best predicts later improvement: conditioning gap,
ensemble disagreement, local coverage/recency, or observed learning progress?
\item Does adaptive allocation improve sample efficiency or weak-region
performance over uniform and compute-matched random allocation?
\item Does any effect arise from rollout preferences, update preferences, or both?
\item How stable are findings across held-out preferences and separately
qualified environment shifts?
\end{enumerate}

\clearpage
\section{Related work and research positioning}

The repository records 22 primary-source full-text readings and structured
mechanism/evidence notes. The following synthesis uses that existing reading
record; it is not a claim that every method has been reproduced locally.

\subsection{Conditioned policies are the starting point, not the novelty}
Dynamic-weight value learning~\cite{abels2019dynamic-weights} and Envelope
Q-learning~\cite{yang2019envelope} establish important antecedents for learning
across linear preference weights. PCN instead conditions a policy on desired
returns and horizon~\cite{reymond2022pcn}; it should not be confused with this
project's preference-conditioned Q-network. CAPQL~\cite{lu2023capql} and
PD-MORL~\cite{basaklar2023pdmorl} further establish modern conditioned continuous-control
methods. PCRL/PreCo addresses objective/preference gradient conflicts through
multi-objective optimization~\cite{yang2025pcrl}. These works rule out a broad
claim that conditioning one policy on trade-offs is new.

More recent mechanisms also constrain interpretation. Hindsight Preference
Replay relabels experience using achieved returns~\cite{shianifar2026hpr},
whereas MORL-FB changes representation learning and preference-guided latent
exploration under a fixed preference distribution~\cite{chen2026morl-fb}.
D$^3$PO uses objective-specific clipped surrogates and a diversity regularizer
to resist collapse~\cite{ambadkar2026d3po}. HPR and D$^3$PO are cited through
author-posted preprints; their status is not silently upgraded. Replay reuse,
optimizer quality, and controllability are potential explanations for an
apparent gain and must be distinguished from allocation itself.

\subsection{Adaptive allocation already has direct antecedents}
PG-MORL predicts improvement to guide population-based training~\cite{xu2020pgmorl}.
GPI prioritization explicitly directs exploration across scalarization
weights~\cite{alegre2023gpi-prioritization}. C-MORL extends sparse regions of a
front through constrained search~\cite{liu2025cmorl}; MPFT allocates further
multi-policy tracking to sparse front regions~\cite{zhao2026mpft}.
PBMORL learns human utility and uncertainty and reallocates multi-policy
optimization toward a region of interest~\cite{hu2026pbmorl}.
Thus, neither adaptive effort over weights nor coverage-directed training is
a defensible stand-alone novelty claim. The narrower proposed study compares
outcome-derived signals inside one shared conditioned learner, separates rollout
from update allocation, and accounts for diagnostic evaluation cost.

\subsection{Curriculum signals need predictive validation}
Value-disagreement curricula motivate uncertainty-driven selection using
multiple value estimators~\cite{zhang2020value-disagreement}. The present shared-trunk
masked heads are correlated; disagreement is a measured score, not a calibrated
epistemic uncertainty guarantee. ALP-GMM~\cite{portelas2020alp-gmm} and
CURIOUS~\cite{colas2019curious} motivate absolute learning/competence progress
as a selection signal. Absolute changes may also capture forgetting or noise,
so predictive association with subsequent improvement must be tested rather
than assumed. Prioritized Level Replay prioritizes environment levels with
learning-potential proxies and staleness~\cite{jiang2021plr}; it is not the
transition-level prioritized replay used in remediation v2.

\subsection{Evaluation must separate the claims}
Environment-context generalization is distinct from interpolation across
preferences~\cite{teoh2025morl-generalization}. A controllability study shows
why front-level metrics alone need not establish that a policy follows its
input weights~\cite{delasheras2026controllability}; this source is a preprint.
Theory for predictable adversarial preferences and preference-free exploration
uses different assumptions from deep-control curriculum allocation and raises
explicit hypervolume semantics~\cite{hu2026predictable-preferences}. Reactive
step-wise weight adversaries define another robustness game, not a physical
dynamics shift~\cite{hu2026reactive-preferences}. These readings inform boundaries;
they do not authorize changing frozen metrics or claiming untested robustness.

\section{Implemented method and experimental controls}

\begin{figure}[htbp]
\centering
\begin{tikzpicture}[node distance=10mm and 9mm]
\node[block] (pref) {Preference sampler\\$w\sim p_t(w)$};
\node[block,right=of pref] (env) {Rollout + environment\\vector reward $r$};
\node[block,right=of env] (replay) {Replay buffer\\transitions + provenance};
\node[block,below=18mm of replay] (q) {Masked-head PCQN\\vector Double-DQN update};
\node[block,left=of q] (scores) {Curriculum diagnostics\\gap / disagreement\\coverage / progress};
\node[block,left=of scores] (alloc) {Budgeted allocation\\uniform-floor mixture};
\draw[flow] (pref)--(env); \draw[flow] (env)--(replay);
\draw[flow] (replay)--(q); \draw[flow] (q)--(scores);
\draw[flow] (scores)--(alloc); \draw[flow] (alloc)--(pref);
\node[below=11mm of scores,align=center,font=\sffamily\small,text=Muted]
{Separate validation selects checkpoints; final-test remains locked.\\
Controls share the preference pool and diagnostic budget.};
\end{tikzpicture}

\caption{Implemented allocation loop at a conceptual level. Allocation can act
on rollout preferences, update preferences, or both; the diagram shows the joint
case. Diagnostic evaluation is budgeted separately from training.}
\end{figure}

\subsection{Preference-conditioned value learning}
The network predicts a vector action value $Q_{\theta,k}(s,a,w)$ for each
masked ensemble head $k$. Greedy evaluation chooses
\[
 a^*(s,w)=\arg\max_a w^\top\frac{1}{K}\sum_{k=1}^{K}Q_{\theta,k}(s,a,w).
\]
Training uses per-head Double-DQN targets and a masked Huber loss. Heads
share a feature trunk; they are not described as fully independent bootstrap
models. Replay preserves behavior preferences, vector rewards, bootstrap masks,
and distinct termination/truncation flags. Time-limit truncation ends a rollout
but does not itself remove Bellman continuation.

\subsection{Allocation signals and controls}
\begin{tabularx}{\textwidth}{@{}p{36mm}Y@{}}
\toprule
\textbf{Signal} & \textbf{Operational interpretation} \\
\midrule
Conditioning gap & Difference between archive-supported utility and the
current policy's measured utility for a preference. \\
Disagreement & Variance of scalarized values across masked heads; calibration
is not assumed. \\
Coverage/recency & Local preference visitation density and time since sampling. \\
Learning progress & Absolute change in measured utility between curriculum checkpoints. \\
\bottomrule
\end{tabularx}

Normalized signals drive temperature-based sampling mixed with a uniform floor.
The implementation separates \texttt{rollout\_only}, \texttt{update\_only}, and
\texttt{rollout\_and\_update}. Uniform-PCQN samples from the same finite
protocol-seeded preference pool. The random curriculum control performs matched
diagnostic evaluations but allocates from random scores. These controls prevent
extra evaluation or a different preference pool from masquerading as an
adaptive-allocation effect.

\subsection{Evaluation and reproducibility}
Training interactions, gradient updates, curriculum-evaluation interactions,
and validation/final-evaluation interactions are counted separately. Pilot-only
calibration froze scales and hypervolume reference points before final testing.
Validation selects checkpoints; the final-test split stays locked until the
complete study and selection gates pass.

Primary measures include hypervolume, expected utility, tail regret, and
learning-curve summaries. Raw and discounted vector returns are retained.
Analysis is paired at training-seed level, emphasizing effect sizes and
confidence intervals rather than treating five seeds as strong significance
evidence. Preference generalization is distinct from environment-shift robustness.

\section{Completed work and verified findings}

\subsection{Progress ledger}
\begin{figure}[htbp]
\centering
\begin{tikzpicture}[node distance=4mm,
stage/.style={draw=Teal,fill=Teal!7,rounded corners,text width=140mm,align=left,inner sep=3mm,font=\sffamily\small}]
\node[stage] (a) {\textbf{Foundation + implementation}\quad Question, 22-source map, protocols, paired controls and provenance.};
\node[stage,below=of a] (b) {\textbf{Primary block}\quad 30 jobs complete; Minecart ore is zero in 14,400 validation episodes.};
\node[stage,draw=Accent,fill=Accent!7,below=of b] (c) {\textbf{Qualification v1}\quad One of three selected seeds passes; all-seeds gate fails.};
\node[stage,draw=Accent,fill=Accent!7,below=of c] (d) {\textbf{Qualification v2 (PER)}\quad Zero of three selected seeds passes; source audit passes 30/30.};
\node[stage,draw=Muted,fill=white,dashed,below=of d] (e) {\textbf{Current stop}\quad Evidence review and versioned scope decision before new training. Final-test locked.};
\foreach \x/\y in {a/b,b/c,c/d,d/e} \draw[flow] (\x)--(\y);
\end{tikzpicture}

\caption{Evidence progression and the present scientific stop. Completed jobs
are not equivalent to a passed behavioral qualification or a supported claim.}
\end{figure}
\begin{tabularx}{\textwidth}{@{}p{39mm}Y@{}}
\toprule
\textbf{Workstream} & \textbf{Evidence-bounded status} \\
\midrule
Research foundation & Revised executable question, architecture contracts,
claim registry, and primary-source literature notes completed. \\
Calibration and plumbing & Both course protocols frozen from non-test
sources; Envelope/GPI-LS and enabled LunarLander shifts passed bounded qualification. \\
Primary block & 30 full-budget jobs completed; source audit and
300-checkpoint validation re-evaluation completed. \\
Scientific comparison & Blocked: all 15 Minecart primary runs had zero ore
at every evaluated checkpoint; zero checkpoints locked. \\
Backbone remediation & V1 and v2 completed and failed their all-seeds rules.
No allocation-effect claim was promoted. \\
Demo and manuscript & Native-sprite checkpoint replays prepared for LunarLander
and Minecart. This document reports progress, not a completed comparative paper
or a qualified Minecart policy. \\
\bottomrule
\end{tabularx}

\subsection{Primary-block failure and localization}
The primary block comprised Uniform-PCQN, matched-random curriculum, and the
full adaptive rule across two environments and five paired seeds. All 30 jobs
completed. Their source audit passed, but both Minecart ore objectives were
constant zero across 14,400 validation episodes spanning all 150 Minecart
checkpoints. The LunarLander block was non-degenerate; the atomic selection
gate nevertheless prohibited partial checkpoint locking.

All 15 diagnostically selected Minecart training archives contained positive
ore returns, while greedy validation did not reproduce them. Of the 14,400
validation episodes, 13,308 (92.42\%) reached the 1,000-step limit. A separate
representative head/action probe made 345,417 action decisions without a
\texttt{Mine} action. These observations are consistent with an underqualified
sparse-credit backbone; they do not identify a treatment effect of allocation.

\subsection{Remediation v1: transient success was insufficient}
V1 ran three fresh Uniform-PCQN seeds for 1,000,000 training interactions each.
All ten checkpoints per seed were evaluated under the predeclared nominal
validation protocol. The selected checkpoints produced:

\begin{center}
\begin{tabular}{@{}rrrrl@{}}
\toprule
\textbf{Seed} & \textbf{Selected step} & \textbf{Ore 0 positive} &
\textbf{Ore 1 positive} & \textbf{Gate} \\
\midrule
20264001 & 800,000 & 0/160 & 0/160 & Fail \\
20264002 & 600,000 & 0/160 & 0/160 & Fail \\
20264003 & 600,000 & 28/160 & 28/160 & Pass \\
\bottomrule
\end{tabular}
\end{center}

V1 failed because two seeds failed. A predeclared probe (seven checkpoints,
1,120 episodes) found 35 cargo-acquisition and 27 delivery episodes at seed
20264003, step 600,000, followed by zero \texttt{Mine} actions at later checkpoints.
This motivated the retention-oriented v2 intervention.

\subsection{Remediation v2: prioritized replay did not qualify the backbone}
V2 changed only uniform replay to proportional prioritized replay, with
$\alpha=0.6$, $\beta:0.4\rightarrow1.0$, and priority epsilon $10^{-6}$.
The masked Huber loss and all other predeclared v1 training/evaluation choices
were preserved. A loss-wording inconsistency was corrected and recorded before
execution; unit importance weights reproduce the v1 loss.

All three jobs completed 1,000,000 training interactions and 995,001 gradient
updates each. The source audit passed all 30 checkpoints. The qualification
evaluated 4,800 nominal-validation episodes, totaling 4,710,905 environment
steps. At the protocol-selected checkpoints:

\begin{center}
\begin{tabular}{@{}rrrrl@{}}
\toprule
\textbf{Seed} & \textbf{Selected step} & \textbf{Ore 0 positive} &
\textbf{Ore 1 positive} & \textbf{Gate} \\
\midrule
20264011 & 1,000,000 & 0/160 & 0/160 & Fail \\
20264012 & 1,000,000 & 0/160 & 0/160 & Fail \\
20264013 & 1,000,000 & 0/160 & 0/160 & Fail \\
\bottomrule
\end{tabular}
\end{center}

Each selected ore objective also had zero distinct positive values. The rule
required at least 8/160 positive episodes \emph{and} two distinct positive
values per ore objective, for every seed. V2 therefore failed. Nonzero
hypervolume alone cannot override this behavioral gate. A first interrupted
attempt of seed 20264011 is preserved separately and must remain disclosed as
additional expenditure; 3,000,000 interactions describes the three completed
runs, not the entire tuning history.

\textbf{Interpretation.} This particular predeclared PER intervention did not
meet the Minecart backbone qualification. The result does not establish that
PER never works, that all MORL methods fail, or that adaptive allocation has
no benefit. It prevents a valid allocation comparison under the current scope.

\begin{figure}[htbp]
\centering
\begin{tikzpicture}
\begin{axis}[width=0.94\textwidth,height=65mm,ymin=0,ymax=34,
ylabel={Positive episodes per ore objective / 160},
symbolic x coords={v1-01,v1-02,v1-03,v2-11,v2-12,v2-13},
xtick=data,ytick={0,8,16,24,32},ymajorgrids,grid style={gray!15},
axis line style={gray!50},tick label style={font=\small},label style={font=\small},clip=false]
\addplot[ycomb,very thick,Teal,mark=*,mark size=3pt,nodes near coords,
every node near coord/.append style={font=\small,above=2pt}]
coordinates {(v1-01,0) (v1-02,0) (v1-03,28) (v2-11,0) (v2-12,0) (v2-13,0)};
\addplot[Accent,dashed,thick] coordinates {(v1-01,8) (v2-13,8)};
\node[anchor=south east,font=\small,text=Accent] at (axis cs:v2-13,8) {Required count: 8};
\end{axis}
\end{tikzpicture}

\caption{Selected-checkpoint ore-positive counts from qualification artifacts.
Both ore objectives have identical counts, so each marker represents either
objective, not their sum. Labels abbreviate seeds 20264001--03 and 20264011--13.
The count threshold is only one gate: two distinct positive values per objective
and all-seeds success are also required. The two batches use different seeds;
this is not a paired estimate of PER's effect.}
\end{figure}

\clearpage
\section{Next steps, risks, and presentation demo}

\subsection{Immediate course deliverable}
For the 19 October checkpoint, the deliverable is this evidence-bounded
progress report: implemented functionality, audited experiments, negative
findings, blockers, and a concrete plan. A finished conference paper or a
successful adaptive-policy demonstration is not claimed.

The frozen failure rule stops current-scope algorithmic remediation. A
supervisor-approved, versioned scope decision is required before v3, a changed
Minecart study, or any additional training batch. Matrix v1 remains permanently
blocked. Checkpoints are not locked, and final-test data remain inaccessible.
The next scientific recommendation must be based on existing evidence rather
than unreviewed tuning until a desired result appears.

\subsection{Native-sprite checkpoint replay demo}
The companion videos render actual MO-Gymnasium environments using saved
checkpoints: LunarLander at 50k/500k interactions (adaptive seed 20262001) and
Minecart at 100k/1M (v2 seed 20264011). Each pair uses one predeclared nominal
validation seed and the existing validation preference nearest the simplex
centre. Both panels use greedy ensemble-mean inference; neither performs
training. These are new replays of saved training stages, not original training
footage or a MORL-versus-non-MORL comparison.

In the fixed LunarLander example, the landing component changes from $-100$
to $+100$; this single trajectory pair is not a success-rate estimate. Both
Minecart examples reach 1,000 steps with zero ore. Actions, cumulative raw
vector returns and terminal states remain visible; failed behavior is retained.
The two MP4s and local video player work without research dependencies. No
training, final-test evaluation or paid review is triggered by opening them.
The predeclared selection and full traces are recorded separately from
qualification evidence and do not unlock any scientific gate.

\subsection{AI-assisted workflow and risk controls}
ARIS skills and a tools-off Codex-to-Claude review bridge have been installed.
Offline checks cover cost/model metadata, transactional reservations, packet
boundaries, quota freshness, interruption cleanup, and resume behavior. These
application controls are not an operating-system isolation guarantee.

Exactly one subsequently authorized supervised batch reviewed four existing
research artifacts. The actual reviewer was \texttt{claude-sonnet-5-5}; CLI
list-basis cost was \$0.0954889, below the \$0.25 invocation cap. The reviewer
found packet reporting internally consistent with caveats, while explicitly
not re-auditing omitted raw tensors or source-audit contents. A historical gate
field still described v1 as predeclared; that wording does not change the latest
failed v2 decision. The coordinator verified that the experiment registry already
records the negative outcome. No claim was promoted.

Historical reservations of \$0.50 were preserved and the one new invocation
reserved \$0.25, totaling \$0.75 within the nominal \$1 envelope. The historical
Console total was user-reported as \$0.24 at cent precision; individual historical
costs were not invented. Completed-state resume retained exactly one new call.
The single-batch authority is consumed and paid dispatch is disabled again.
Auto-reload remains off according to the user's Console evidence. No recurring
scheduler is enabled. Unattended operation is not qualified: OS isolation for
process-group escape or simultaneous guardian termination remains unverified,
and one supervised review does not certify every research workflow.

\subsection{Decisions to request from the supervisor}
The team should ask whether the course checkpoint should present the blocked
negative result under the frozen scope, or whether a separately versioned scope
amendment is desired. A further clarification is whether a presentation of the
implemented pipeline and diagnostic evidence meets the initial-demo requirement.
These questions seek direction; they do not implement a scope change or authorize
v3, new training, checkpoint locking, or final-test access.

\section{Conclusion}
The project has progressed from a proposal to an executable, provenance-aware
research system with paired controls, frozen evaluation choices, tested
implementation invariants, and reproducible negative evidence. The main
blocker is scientific rather than a missing training script: the Minecart
backbone has failed two predeclared qualifications. Reporting that failure and
preserving the data are necessary progress toward a trustworthy study.
The course checkpoint can demonstrate this engineering and diagnostic work
without claiming an adaptive-allocation improvement that the evidence does
not establish.

\section*{Evidence and reproducibility notes}
\addcontentsline{toc}{section}{Evidence and reproducibility notes}
Training source: \texttt{25328e8}. The companion \texttt{EVIDENCE\_INDEX.md}
links primary validation, checkpoint-lock decisions, both qualifications,
behavior diagnostics, source audit and raw returns. This report changes no
protocol, completes no research phase and unlocks no final-test data.

\clearpage
